% !TeX root = ../main.tex
\section{Related Work}

\subsection{Underwater Image Enhancement}
Early UIE research predominantly relied on physics-inspired modeling and hand-crafted priors. Based on the Jaffe--McGlamery image formation model, many methods estimate background light and medium transmission to invert attenuation and backscatter effects \cite{r12,r4,r14}, while non-physical approaches improve perceptual quality by manipulating intensity and color statistics such as histogram equalization and Retinex-style illumination correction \cite{r15,r17}. With public benchmarks and improved training recipes, deep learning has become the dominant paradigm \cite{r18,r6}, and representative CNN and Transformer designs improve detail recovery, global color consistency, and contrast \cite{r0,r7}. Hybrid methods further inject medium-dependent cues such as transmission or multi-color embeddings to enhance robustness \cite{r42,r40,r41}, but UIE models still suffer from local over-enhancement, unstable color correction, and performance degradation under unseen conditions.

\subsection{Diffusion Models}
Diffusion models have achieved strong performance in image synthesis and restoration by iterative denoising with powerful learned priors \cite{r20,r52,r21}, and their flexibility has enabled diverse conditional generation frameworks for low-level vision \cite{r22}. In underwater enhancement, recent studies use conditional diffusion to improve texture realism and color restoration, sometimes with physical priors or Transformer backbones for better robustness \cite{r9,r24,r8,r10}. The main limitation is efficiency, because conventional diffusion requires many sequential sampling steps. Existing acceleration strategies include improved solvers \cite{r43,r44}, latent diffusion \cite{r47}, and distillation or consistency training for one- or few-step inference \cite{r45,r46,r48}, yet deployment on resource-constrained underwater platforms remains difficult due to model size, iterative computation, and unstable conditioning across diverse water environments.

\subsection{Rectified Flow}
Rectified flow provides an efficient alternative by learning deterministic transport from noise to data through a regression-style objective, enabling fast ODE sampling with fewer function evaluations than diffusion \cite{r11}. Related formulations such as stochastic interpolants and flow matching further improve stability and controllability \cite{r28,r29}, and recent results in InstaFlow and FlowIE show that one-step or few-step generation can retain strong perceptual quality in practical enhancement pipelines \cite{r30,r49}. Although low-step performance continues to improve through better training and parameterization strategies \cite{r50,r51}, underwater enhancement is still challenging because of limited real paired supervision, large water-type variation, and unreliable physical cue estimation; these factors can amplify artifacts when efficient generative models are conditioned incorrectly. This motivates our design that combines rectified flow with physics-guided multi-condition inputs and adaptive masking for more stable restoration.
